\section{What is open}\label{sec:open}

\begin{enumerate}[leftmargin=1.6em,itemsep=0.4em]
\item \emph{An off-lattice terminal spot.} By Remark~\ref{rem:bary} the chart
represents every $s\in[K^0,K^{N+1}]$ at expiry, and the state that does so is
$\varpi(s)$, carried by the two nodes straddling $s$. Whether a front volatility
exists driving $q^1$ to $\varpi(s)$ for $s$ strictly inside a cell, rather than to
a single node as \eqref{eq:nufront} does, is open. The obstruction is that
exact settlement forces $q^1_{T_1}=\varpi(S_{T_1})$ and hence
$\int_0^{T_1}(\eta^{1,i})^2\dd t<\infty$ for both surviving nodes, which by the
argument of Proposition~\ref{prop:collapse} they cannot both satisfy when the
volatility is proportional to distance from the spot. Until it is answered the
model of Proposition~\ref{prop:collapse} pins, as \cite{Wissel2007} does, and
his question of a front that does not pin the spot remains open.

\item \emph{The continuum limit.} The constant $(M-1)(N+1)$ grows with the grid. Whether $\sup_N\Theta_j<\infty$ on a class of
realistic surfaces, which is what a continuum statement would need, is open.
\end{enumerate}
